\documentclass[11pt]{article}
\usepackage{fontspec}
\setmainfont{DejaVu Serif}[Scale=0.92]
\setsansfont{DejaVu Sans}[Scale=0.90]
\setmonofont{DejaVu Sans Mono}[Scale=0.84]
\usepackage[a4paper,margin=2.4cm]{geometry}
\usepackage{xcolor}
\definecolor{ink}{HTML}{1A1A1A}
\definecolor{rule}{HTML}{C8C8C8}
\definecolor{ansbg}{HTML}{F4F7FB}
\definecolor{ansbd}{HTML}{3A6EA5}
\definecolor{vpbg}{HTML}{FAFAF7}
\definecolor{vpbd}{HTML}{BDBDB2}
\usepackage[most]{tcolorbox}
\usepackage{longtable,booktabs,array}
\usepackage{enumitem}
\usepackage{microtype}
\usepackage{titlesec}
\usepackage{parskip}
\usepackage{tocloft}
\setlength{\cftsecnumwidth}{2.6em}
\renewcommand{\cftsecfont}{\normalfont}
\renewcommand{\cftsecpagefont}{\normalfont}
\usepackage[hidelinks]{hyperref}
\titleformat{\section}{\sffamily\large\bfseries\color{ink}}{\thesection.}{0.6em}{}
\titlespacing*{\section}{0pt}{2.2ex plus 1ex minus .2ex}{1.2ex}
\titleformat{\subsection}{\sffamily\normalsize\bfseries\color{ink}}{}{0em}{}
\titlespacing*{\subsection}{0pt}{1.6ex plus .5ex}{0.7ex}
\newtcolorbox{answerbox}{colback=ansbg,colframe=ansbd,boxrule=0.4pt,left=8pt,right=8pt,top=6pt,bottom=6pt,
  arc=1pt,before skip=6pt,after skip=8pt,fontupper=\normalsize}
\newtcolorbox{abstractbox}{colback=white,colframe=rule,boxrule=0.3pt,left=8pt,right=8pt,top=6pt,bottom=6pt,
  arc=0pt,before skip=4pt,after skip=8pt,fontupper=\small\itshape}
\newtcolorbox{vpbox}{colback=vpbg,colframe=vpbd,boxrule=0.4pt,left=8pt,right=8pt,top=6pt,bottom=6pt,
  arc=1pt,before skip=6pt,after skip=8pt,fontupper=\small}
\setlength{\emergencystretch}{3em}
\hyphenpenalty=1500
\sloppy
\title{\sffamily\bfseries Aging and Senescence:\\[2pt]
  {\large\mdseries the Systemic Decline of Homeostatic Setpoints over the Lifespan}}
\author{Young Jae Lee \\ \small ORCID 0009-0002-7535-8245}
\date{\small Version 1.4.0 \;\textbullet\; DOI 10.5281/zenodo.20756155 \;\textbullet\; CC BY 4.0 \;\textbullet\; 2026}
\begin{document}
\maketitle
\thispagestyle{empty}
\begin{abstract}\noindent
Aging is read on the VP substrate as the systemic, monotonic decline of $\gamma$-defended homeostatic
setpoints. From sequence the promoter $\gamma$ of the aging masters (TP53, CDKN2A, FOXO3, TERT) is measured
and fixes each node's barrier $\gamma^2/4$ and reservoir dwell $\sim\gamma^{1.5}$; the dynamics --- setpoint
drift, the stuck senescence attractor, finite-reservoir depletion, a convex age risk-multiplier --- are
deterministic simulations seeded at 19. A cross-species discriminant finds human aging genes are not special
and the longevity switch is TP53 copy number, not promoter $\gamma$; an observation-only archaic$\leftrightarrow$
present-day comparison finds the masters read the same $\gamma$ in dated genomes; and a telomere deep-dive
locates the telomere as the keystone of aging \emph{dynamics}, not $\gamma$. The closing section draws the
epistemic conclusion: the genome fixes the ruler but not the realized lifespan, which is underdetermined by
the $\gamma$ axis (the lever lives off-axis in dosage and reservoir dynamics) rather than merely uncomputed.
Everything is reproduced bit-for-bit (SEED=19, $2\times$sha256 identical) and graded honestly
([V] verified / [L] cited / [O] open-with-obstacle).
\end{abstract}
\vspace{1em}
\tableofcontents
\clearpage


\section{Aging as homeostatic setpoint decline}

\begin{answerbox}Aging on the jamming substrate is the slow loss of defense gain of every homeostatic setpoint, plus the accumulation of cells stuck in irreversible R19 attractors (senescence). It is the dominant risk multiplier for every pathology kernel. No new organs are emerged here: node identities are inherited from DNA, and this capstone layer adds the temporal dynamics.\end{answerbox}

\begin{abstractbox}The R19 field ṡ = γs − s³ + h has two basins separated by a barrier γ²/4; a defended homeostatic setpoint is the healthy basin, and “defense gain” is carried by γ. Aging is modeled as a slow decline of that gain, which both lets the defended value drift and lowers the catastrophe threshold spinodal 2(γ/3)\textsuperscript{1.5}.\end{abstractbox}

\subsection*{What aging is on this substrate}

Aging is not an organ. It is a process that runs on the same vacuum-jamming R19 substrate used across the framework, where every switch obeys the normal-form field ṡ = γs − s³ + h. That field has two stable basins; the energy barrier between them is γ²/4, and the drive |h| past which one basin vanishes is the spinodal 2(γ/3)\textsuperscript{1.5}.

A defended homeostatic setpoint — glucose, pressure, calcium, temperature — is the healthy basin held against perturbation. Its restoring strength (loop gain) is carried by γ. Aging is the slow decline of that gain over the lifespan.

\begin{vpbox}\textbf{barrier(γ) = γ\textsuperscript{2}/4} — the energy barrier between the two basins; it sets how stable a defended setpoint is. [F] forced substrate, §1\end{vpbox}

\begin{vpbox}\textbf{spinodal(γ) = 2(γ/3)\textsuperscript{1.5}} — the drive |h| past which one R19 basin disappears, so the flip is discontinuous (a saddle-node). [F] forced substrate, §1\end{vpbox}

\subsection*{Two failure modes from one decline}

A falling gain does two things at once. The barrier shrinks, so the defended value drifts inside its basin — the gradual, multi-system signature of aging (§3). And the spinodal shrinks, so a chronic stressor that was survivable when young can flip the setpoint discontinuously into the pathological basin — the sudden failure event (§3, §12).

Two further consequences define the program: cells can cross into an irreversible arrested basin and accumulate (senescence, §4), and finite reservoirs deplete to a replicative limit (§5).

\subsection*{Nodes and seams}

This layer watches four measured masters and one diffuse node. The masters are TP53 (cellular senescence), CDKN2A (the arrest switch), TERT (telomere maintenance), and FOXO3 (longevity signaling); their identity γ is owned by DNA and never fitted (§2). The diffuse node is the systemic drift of every imported setpoint.

\begin{vpbox}\textbf{γ (promoter identity)} — −mean nearest-neighbour stacking ΔG37 (SantaLucia 1998) over the promoter window; the DNA Layer-1 identity that fixes each node's barrier and dwell. [V] verified DNA Layer 1 (DOI)\end{vpbox}

The single seam this package is the source of truth for is the aging risk multiplier: it raises the crossing rate of every oncology and pathology kernel across the framework (§7).

\section{Node atlas: measured promoter gamma}

\begin{answerbox}The four aging masters carry measured promoter γ: TP53 1.4298, CDKN2A 1.4424, TERT 1.5539, FOXO3 1.5942 — each computed as −mean SantaLucia nearest-neighbour ΔG37 over the TSS−2000..+500 window and never fitted. Emergence order by ascending γ is cellular senescence, arrest switch, telomere maintenance, longevity signaling.\end{answerbox}

\begin{abstractbox}γ is the promoter stacking-energy identity from DNA Layer 1; it fixes each node's barrier γ²/4 and dwell γ\textsuperscript{1.5}. The four values are vendored from NCBI RefSeq accessions and reproduce bit-for-bit from the committed promoter cache.\end{abstractbox}

\subsection*{The masters and their identity}

Each aging master is a real promoter, and its γ is a measurement, not a parameter. The window is fixed (TSS−2000 to +500, 2501 bp), the nearest-neighbour table is SantaLucia 1998, and the value is −mean ΔG37 over the window. The same pipeline produces the DNA gene-clock γ.

\begin{center}\footnotesize
\begin{longtable}{@{}>{\raggedright\arraybackslash}p{0.158\linewidth} >{\raggedright\arraybackslash}p{0.158\linewidth} >{\raggedright\arraybackslash}p{0.158\linewidth} >{\raggedright\arraybackslash}p{0.158\linewidth} >{\raggedright\arraybackslash}p{0.158\linewidth} >{\raggedright\arraybackslash}p{0.158\linewidth}@{}}
\toprule
\textbf{Gene} & \textbf{Node} & \textbf{γ} & \textbf{Accession} & \textbf{Strand} & \textbf{Dynamics} \\
\midrule
TP53 & cellular senescence & 1.429832 & NC\_000017.11 & - & stuck-attractor \\
CDKN2A & senescence arrest switch & 1.442444 & NC\_000009.12 & - & stuck-attractor \\
TERT & telomere maintenance & 1.553876 & NC\_000005.10 & - & reservoir-depletion \\
FOXO3 & longevity signaling & 1.594156 & NC\_000006.12 & + & maintenance \\
\bottomrule
\end{longtable}
\end{center}

\begin{vpbox}\textbf{γ (promoter identity)} — −mean nearest-neighbour stacking ΔG37 (SantaLucia 1998) over the promoter window; the DNA Layer-1 identity that fixes each node's barrier and dwell. [V] verified DNA Layer 1 (DOI)\end{vpbox}

\subsection*{Why the order matters}

γ sets two derived quantities. The barrier γ²/4 fixes how hard it is to leave a basin (state stability), and the dwell γ\textsuperscript{1.5} fixes how long a switch runs (reservoir capacity, §5). Sorting the masters by γ gives a fixed emergence order: cellular senescence (TP53) and the arrest switch (CDKN2A) sit lowest, then telomere maintenance (TERT), then longevity signaling (FOXO3).

\begin{center}\ttfamily\small γ\textsubscript{TP53} = 1.429832  |  barrier = γ²/4 = 0.5111\end{center}

\subsection*{The diffuse node}

The fifth node, systemic setpoint drift, has no single master gene. It is a circuit-level, derived quantity: the aggregate decline of every defended setpoint imported from the other packages. It is the object that §3 and §8 quantify.

\section{Setpoint drift: the unifying signature}

\begin{answerbox}Under a slow loss of loop gain, a defended setpoint first creeps inside its basin and then flips catastrophically once the gain falls past the spinodal. One law reproduces both the gradual multi-system drift of aging and the discontinuous failure event. The measured flip jump is 0.79; the mechanism is [V], the absolute rate is [O].\end{answerbox}

\begin{abstractbox}Sweeping a falling gain through the R19 field gives a monotone creep of the defended state while the barrier γ²/4 shrinks, followed by a discontinuous jump of 0.790 when the spinodal is crossed. The same decline thus yields both the slow drift of aging and its sudden onset events — one signature, two faces.\end{abstractbox}

\subsection*{Creep, then catastrophe}

The defended value does not fall smoothly to failure. While the basin still exists, the setpoint creeps — a slow, monotone drift that reads as ordinary aging across many systems at once. This is the unifying signature: one declining gain, expressed everywhere a setpoint is defended.

Then the basin disappears. Once the gain drops past the spinodal, the healthy basin is annihilated in a saddle-node event and the state jumps discontinuously to the pathological basin. In the run this jump is 0.7904 — a sudden failure on top of the slow drift.

\begin{vpbox}\textbf{spinodal(γ) = 2(γ/3)\textsuperscript{1.5}} — the drive |h| past which one R19 basin disappears, so the flip is discontinuous (a saddle-node). [F] forced substrate, §1\end{vpbox}

\subsection*{Why aging looks gradual but fails suddenly}

This resolves a familiar tension. Function appears to decline smoothly for decades (the creep), yet clinical failure often arrives as an event — a fall, a fracture, a decompensation. Both come from the same monotone loss of gain; the creep is the sub-spinodal regime and the event is the crossing.

The shape is reproduced and graded [V]. The absolute calendar rate of the decline is not derivable here — it needs an external clock — and is graded [O] (see §13).

\section{Senescence as a stuck attractor}

\begin{answerbox}Cellular senescence is a cell that has crossed into an irreversible arrested R19 basin. Past the spinodal the proliferative basin disappears, so the arrest is one-way — a sub-spinodal pulse is still reversible, but a supra-spinodal crossing is not — and arrested cells accumulate monotonically, reaching a fraction of 0.391 in the run.\end{answerbox}

\begin{abstractbox}Driving a cell past the spinodal annihilates the proliferative basin, making the arrested state an absorbing attractor (return is impossible), while a sub-spinodal pulse relaxes back. Repeated exposure therefore accumulates arrested cells over time, with the arrested fraction rising monotonically to 0.391.\end{abstractbox}

\subsection*{Irreversibility is a saddle-node, not a threshold}

Senescence is permanent for a structural reason. Below the spinodal the proliferative and arrested basins coexist, so a perturbation that pushes a cell toward arrest can relax back. Above the spinodal the proliferative basin no longer exists, so the cell cannot return — the arrested state is absorbing.

\begin{vpbox}\textbf{spinodal(γ) = 2(γ/3)\textsuperscript{1.5}} — the drive |h| past which one R19 basin disappears, so the flip is discontinuous (a saddle-node). [F] forced substrate, §1\end{vpbox}

The run confirms both halves: a supra-spinodal drive is one-way, and a sub-spinodal pulse is reversible. Irreversibility is therefore not an assumption; it follows from which basins exist at a given drive.

\subsection*{Accumulation and SASP}

Because each crossing is permanent, arrested cells accumulate. The arrested fraction rises monotonically with continued exposure, reaching 0.391 in the simulation. The senescence-associated secretory phenotype then raises neighbours' crossing hazard, coupling senescence to the risk multiplier of §7.

The irreversibility and the monotone accumulation are graded [V]; the absolute per-year rate is [O].

\section{Reservoir depletion: the telomere clock}

\begin{answerbox}Each reservoir — telomere/TERT, senescence/CDKN2A, apoptosis/TP53, longevity/FOXO3 — has finite capacity proportional to γ\textsuperscript{1.5}. Deeper-γ wells hold more and last longer, but every reservoir depletes to zero. TERT (γ 1.5539) lasts 89 steps and TP53 (γ 1.4298) lasts 78; telomere attrition is reservoir exhaustion, not a privileged clock.\end{answerbox}

\begin{abstractbox}The dwell γ\textsuperscript{1.5} sets a finite reservoir capacity for each node. Depletion time increases monotonically with γ, yet all four reservoirs reach zero, so the replicative limit is a generic consequence of a finite well rather than a dedicated counting mechanism.\end{abstractbox}

\subsection*{Finite wells}

A switch that runs for a dwell γ\textsuperscript{1.5} draws on a finite reservoir. The deeper the well (larger γ), the more it holds and the longer it lasts — but the reservoir is finite, so it empties. This is the substrate's version of the replicative limit.

\begin{vpbox}\textbf{dwell(γ) ≈ γ\textsuperscript{1.5}} — how long a switch runs before its finite reservoir empties; it sets reservoir capacity (the replicative limit). [F] forced DNA Layer 1 (DOI)\end{vpbox}

\begin{center}\footnotesize
\begin{longtable}{@{}>{\raggedright\arraybackslash}p{0.237\linewidth} >{\raggedright\arraybackslash}p{0.237\linewidth} >{\raggedright\arraybackslash}p{0.237\linewidth} >{\raggedright\arraybackslash}p{0.237\linewidth}@{}}
\toprule
\textbf{Reservoir} & \textbf{γ} & \textbf{Capacity} & \textbf{Steps to empty} \\
\midrule
apoptosis (TP53) & 1.4298 & 1.554 & 78 \\
senescence (CDKN2A) & 1.4424 & 1.575 & 79 \\
telomere (TERT) & 1.5539 & 1.761 & 89 \\
longevity (FOXO3) & 1.5942 & 1.830 & 92 \\
\bottomrule
\end{longtable}
\end{center}

\subsection*{Telomere attrition is not special}

The telomere/TERT reservoir is the longest-lived of the four because its γ is high, but it depletes by the same rule as the others. Depletion time rises with γ across all four wells. The telomere clock is therefore reservoir exhaustion, not a privileged counting device — consistent with the cross-species result that TERT γ is not a longevity switch (§9).

The finiteness and the γ-ordering are graded [V]; the absolute attrition rate in years is [O].

\section{Hallmarks of aging on the substrate}

\begin{answerbox}All ten hallmarks of aging map onto substrate mechanisms with no orphans. Genomic instability and cellular senescence are R19 mis-flips and stuck attractors (RA2); telomere attrition and stem-cell exhaustion are reservoir depletion (RA3); loss of proteostasis, deregulated nutrient sensing and epigenetic drift are loop-gain loss and setpoint drift (RA1).\end{answerbox}

\begin{abstractbox}Each of the ten classical hallmarks is assigned to a single substrate phenomenon already reproduced in RA1–RA3, leaving zero orphans. The map is structural — every hallmark points to a mechanism that the engine demonstrates, with three hallmarks carrying an [O] tail for their absolute magnitude.\end{abstractbox}

\subsection*{One substrate, ten hallmarks}

The hallmarks of aging are usually a list. Here they are consequences of three substrate behaviours: switch mis-flips and absorbing arrest (RA2), finite-reservoir depletion (RA3), and the slow loss of loop gain that drifts setpoints (RA1). Each hallmark maps to exactly one, with no remainder.

\begin{center}\footnotesize
\begin{longtable}{@{}>{\raggedright\arraybackslash}p{0.237\linewidth} >{\raggedright\arraybackslash}p{0.237\linewidth} >{\raggedright\arraybackslash}p{0.237\linewidth} >{\raggedright\arraybackslash}p{0.237\linewidth}@{}}
\toprule
\textbf{Hallmark} & \textbf{Substrate mechanism} & \textbf{Axis} & \textbf{Grade} \\
\midrule
genomic instability & R19 switch mis-flips (errored basin crossings) & RA2/substrate & [V] \\
telomere attrition & finite reservoir clock, capacity \textasciitilde{} gamma\textasciicircum{}1.5 & RA3 & [V] \\
epigenetic alteration & setpoint of the cis-drive threshold drifts & RA1 & [V] \\
loss of proteostasis & basin/loop-gain degradation (shallower barrier) & RA1 & [V] \\
deregulated nutrient sensing & FOXO3/IGF longevity-loop gain decline & RA1(longevity) & [V] \\
mitochondrial dysfunction & energy-supply gain drop -> setpoint drift & RA1 & [V]/[O] \\
cellular senescence & stuck (irreversible) arrested R19 attractor + SASP & RA2 & [V] \\
stem-cell exhaustion & reservoir depletion to the replicative limit & RA3 & [V] \\
altered intercellular comm. & SASP raises neighbours' crossing hazard (coupling) & RA2/RA5 & [V]/[O] \\
chronic inflammation & immunosenescence: clearance-loop gain decline & RA1/RA5 & [V]/[O] \\
\bottomrule
\end{longtable}
\end{center}

The structural completeness — every hallmark resolves to a demonstrated mechanism — is graded [V]. Where a hallmark's absolute size needs external calibration (mitochondrial output, intercellular coupling strength, inflammatory tone), the mapping carries a [V]/[O] tail.

\section{Aging as the universal risk multiplier}

\begin{answerbox}As barriers shrink with age, the crossing rate of every pathology kernel rises, so cumulative incidence is convex and steepening — a 52.3× rise between ages 40 and 80 in the model, before immunosenescence is added. This is why a single time axis multiplies risk across unrelated diseases. The shape is [V]; absolute incidence is [O].\end{answerbox}

\begin{abstractbox}A shrinking barrier raises the Kramers crossing hazard, and accumulated crossings give a convex age-incidence curve. Cumulative incidence rises from 0.0056 at age 40 to 0.292 at age 80, a 52.3× increase, steepened further by immunosenescence (the seam to the immune package).\end{abstractbox}

\subsection*{Why one axis multiplies many risks}

Most major diseases share a single dominant risk factor: age. On this substrate that is not a coincidence. Every pathology is a barrier-crossing, and aging lowers every barrier, so the crossing hazard of unrelated kernels rises together. The accumulated crossings produce the steep, convex age-incidence curve seen in oncology.

In the model, cumulative incidence climbs from 0.0056 at 40 to 0.292 at 80 — a 52.3× rise — and immunosenescence steepens it further by weakening the clearance loop (the immune seam).

\subsection*{What is claimed}

The convex, steepening shape is reproduced and graded [V]. The absolute incidence magnitude is not derivable in-package — it needs external calibration — and is graded [O] with that obstacle stated (§13). This is the cross-cutting seam the package owns for the rest of the framework.

\section{The rate of aging: one biological age?}

\begin{answerbox}Across systems, aging is dominated by a single shared rate — 0.889 of the variance in the model — with smaller system-specific residuals. A biological-age axis exists and is dominant, but it is not the whole story: organs retain measurable independent rates. The shared-plus-residual structure is [V]; the absolute mapping is [O].\end{answerbox}

\begin{abstractbox}Decomposing per-system aging rates yields a dominant shared component of 0.889 plus residual per-system variation. This supports a single biological-age axis distinct from chronological age, while preserving organ-specific divergence — neither a pure global clock nor fully independent aging.\end{abstractbox}

\subsection*{Shared rate versus per-system rates}

Do all systems age at one rate, or each at its own? The decomposition gives a clear but two-sided answer: a dominant shared rate accounting for 0.889 of the variance, plus real system-specific residuals. A biological-age axis exists and dominates, yet organs are not locked to it.

This matches observation: biological age tracks chronological age strongly but imperfectly, and individual organs can age faster or slower than the body as a whole. The structure — one dominant axis plus residuals — is graded [V]; mapping it to an absolute clock is [O].

\section{Are human aging genes special?}

\begin{answerbox}No. Across ten mammals spanning 3.8 to 122 years, every human aging-gene promoter γ sits inside the mammalian distribution (|z|<1), and there is no discontinuous γ longevity switch. The core gate TP53 is nearly flat (CV 3.5\%; elephant ≈ human ≈ mouse). The real long-lived switch is off the γ axis: TP53 copy number.\end{answerbox}

\begin{abstractbox}Promoter γ for TP53, CDKN2A, FOXO3 and TERT was measured across ten mammals. Human is within one standard deviation on every gene, no clean short/long threshold appears, and TP53 varies by only CV 3.5\% across a 4–211 year span. The elephant's cancer resistance instead comes from \textasciitilde{}20 TP53 copies at essentially unchanged per-copy γ 1.4269.\end{abstractbox}

This chapter is the package's direct answer to the question that motivated the study: are human aging genes different from other animals, or is there a special switch? It is the headline result.

\subsection*{The cross-species table}

The same promoter-γ pipeline was run for the four masters across ten mammals from the brown rat (3.8 yr) to the human (122 yr), with the bowhead whale (211 yr) listed for span. Coverage is honest: bowhead has no annotated target promoters in NCBI Gene, and CDKN2A is unannotated in several species (the free-text hits were a different gene, CDKN2AIP).

\begin{center}\footnotesize
\begin{longtable}{@{}>{\raggedright\arraybackslash}p{0.158\linewidth} >{\raggedright\arraybackslash}p{0.158\linewidth} >{\raggedright\arraybackslash}p{0.158\linewidth} >{\raggedright\arraybackslash}p{0.158\linewidth} >{\raggedright\arraybackslash}p{0.158\linewidth} >{\raggedright\arraybackslash}p{0.158\linewidth}@{}}
\toprule
\textbf{Species} & \textbf{MLSP} & \textbf{TP53} & \textbf{CDKN2A} & \textbf{FOXO3} & \textbf{TERT} \\
\midrule
bowhead whale & 211.0 & – & – & – & – \\
human & 122.5 & 1.4298 & 1.4424 & 1.5942 & 1.5539 \\
African elephant & 65.0 & 1.4269 & – & 1.6129 & 1.5207 \\
chimpanzee & 59.4 & 1.4319 & 1.5539 & 1.5992 & 1.5528 \\
little brown bat & 34.0 & 1.4566 & 1.4275 & 1.5594 & 1.5046 \\
naked mole rat & 31.0 & 1.4053 & – & 1.5971 & 1.5719 \\
dog & 24.0 & 1.5261 & – & 1.7217 & 1.5867 \\
cattle & 20.0 & 1.4253 & 1.4670 & 1.5436 & 1.5530 \\
gray opossum & 5.1 & 1.3212 & 1.3950 & 1.4306 & 1.3778 \\
house mouse & 4.0 & 1.4194 & 1.5100 & 1.5758 & 1.4227 \\
brown rat & 3.8 & 1.4118 & 1.2922 & 1.4512 & 1.3373 \\
\bottomrule
\end{longtable}
\end{center}

\subsection*{Human is not special}

On every gene the human value sits inside the mammalian spread, within one standard deviation, and no single-gene γ–lifespan trend survives a permutation test.

\begin{center}\footnotesize
\begin{longtable}{@{}>{\raggedright\arraybackslash}p{0.136\linewidth} >{\raggedright\arraybackslash}p{0.136\linewidth} >{\raggedright\arraybackslash}p{0.136\linewidth} >{\raggedright\arraybackslash}p{0.136\linewidth} >{\raggedright\arraybackslash}p{0.136\linewidth} >{\raggedright\arraybackslash}p{0.136\linewidth} >{\raggedright\arraybackslash}p{0.136\linewidth}@{}}
\toprule
\textbf{Gene} & \textbf{n} & \textbf{human γ} & \textbf{CV} & \textbf{human z} & \textbf{pctile} & \textbf{perm p} \\
\midrule
TP53 & 10 & 1.4298 & 3.52\% & +0.09 & 60 & 0.114 \\
CDKN2A & 7 & 1.4424 & 5.85\% & +0.01 & 43 & 0.352 \\
FOXO3 & 10 & 1.5942 & 5.27\% & +0.31 & 50 & 0.073 \\
TERT & 10 & 1.5539 & 5.84\% & +0.64 & 70 & 0.126 \\
\bottomrule
\end{longtable}
\end{center}

Three of four genes show overlapping short- and long-lived values (no clean threshold). TP53 in particular is nearly flat across the whole span (CV 3.5\%): elephant, human and mouse share essentially the same core gate. All four masters do, however, lean the \textit{same} way — every per-gene rank correlation with lifespan is positive — so the honest question is not whether one gene switches but whether that shared lean is real or an artifact.

\subsection*{Robustness audit: the shared lean is confound, not signal}

A dedicated audit re-derives γ bit-for-bit and then corrects the confounds the per-gene tests ignored. First, on this panel γ is essentially a GC-content proxy (ρ(γ, GC) = 0.96–1.00 across genes), so the “γ axis” is a GC axis. Second, the four-gene combined lean — a test the original analysis never ran — reaches only a borderline \textit{raw} ρ = 0.64 (permutation p = 0.054), and it does not survive correction: body-mass adjustment (the intrinsic-longevity residual) drops it to p = 0.16, and Felsenstein independent contrasts on a dated mammal tree collapse it to corr = 0.17 (p = 0.55) — the short-lived rodents are a single clade, so the raw trend was phylogenetic pseudoreplication. Out of sample, the four γ values classify long- vs short-lived species at only 80\% leave-one-out accuracy (label-permutation p = 0.11), not above chance. The human-not-special / no-switch conclusion is therefore robust to every confound this small panel lets us test, and the original “no trend” wording is sharpened to “a weak shared lean that is fully confound-attributable.”

One honest caveat survives the audit. Of the four masters, \textbf{TERT (telomere maintenance) is the most lifespan-leaning}: it is the only gene whose lean is \textit{not} removed by body-mass correction (size-corrected ρ = 0.54) and it carries the largest phylogenetic-contrast correlation (0.58). It still never reaches significance (p ≈ 0.12), so it is a lead for a larger panel, not a result — graded [O]. It is consistent with known biology that the telomere lever on lifespan is real but acts off the promoter-γ axis: large mammals suppress somatic telomerase to resist cancer (Gomes 2011), a regulatory and dosage effect rather than a promoter-sequence one — the same off-axis pattern as the TP53 copy-number switch below.

\subsection*{The real switch is off the γ axis}

Where, then, does extreme longevity come from? Not from rewriting the promoter. The elephant carries \textasciitilde{}20 functional TP53 copies at an essentially unchanged per-copy γ of 1.4269 (human 1.4298), giving roughly twenty times the effective gate dosage at a per-copy barrier of 0.5111. The switch is copy number, not promoter γ.

\begin{vpbox}\textbf{γ\textsubscript{TP53} ≈ 1.427 (mammal-invariant)} — the apoptosis/senescence gate barrier is nearly species-invariant across a 4–211 yr lifespan span. [V] verified node atlas, §2\end{vpbox}

The framework's thesis holds cross-species: a conserved identity γ substrate plus divergent dynamics (dosage, loop gain, reservoir size, senescence rate) produces the lifespan spectrum. Human-not-special, TP53 flatness, and the audited finding that the weak shared lean is fully confound-attributable are graded [V]; the residual TERT lean is [O] (n = 10, p ≈ 0.12, awaiting a larger panel); the off-γ-axis longevity levers are [L] — TP53 copy number (Abegglen 2015 JAMA; Sulak 2016 eLife) and somatic telomerase suppression in large mammals (Gomes 2011).

\section{Archaic and present-day aging promoters}

\begin{answerbox}Across a cross-sectional set of seven dated, sequenced individuals — the present-day reference, three archaic Neanderthal genomes, one archaic Denisovan, and two dated modern-human genomes — the four aging-master promoter γ values sit in a very narrow band on every gene (each per-gene range < 0.0016). The senescence/apoptosis gate TP53 and the arrest switch CDKN2A read identically in the two dated modern-human genomes and the present-day reference; TERT carries the most archaic promoter substitutions (13). This is a measured snapshot — the only claims are the measured γ in each individual, its spread per gene, and which positions differ.\end{answerbox}

\begin{abstractbox}The same promoter-γ pipeline (−mean SantaLucia NN ΔG37 over TSS−2000..+500) is measured in each individual, where each promoter is the GRCh37 reference plus that individual's homozygous-derived substitutions. The set is examined as a snapshot; shared genotype states are reported as observed co-occurrences, never as anything beyond the co-occurrence itself.\end{abstractbox}

\begin{vpbox}\textbf{Observation only (constitution)} — this chapter compares a cross-sectional SET of dated individuals. Each is one sequenced genome with one measured promoter γ, reported as a present-state. The package makes no claim about how any state arose or about any process relating the individuals; where two individuals share a genotype state, that is reported as an observed co-occurrence, never as anything beyond the co-occurrence itself. [V] verified\end{vpbox}

\subsection*{The cross-sectional set}

Seven high-coverage genomes carry the four aging masters at a known date: the present-day modern-human reference (GRCh37), three archaic Neanderthal individuals (Altai, Vindija 33.19, Chagyrskaya), one archaic Denisovan, and two dated modern-human genomes (Ust'-Ishim ≈ 45 kya, Loschbour ≈ 8 kya). Each promoter is reconstructed as the GRCh37 reference plus that individual's homozygous-derived substitutions (genotype 1/1, FILTER pass); γ is then measured by the identical pipeline used for the node atlas (§2) and the cross-species panel (§9).

\begin{center}\footnotesize
\begin{longtable}{@{}>{\raggedright\arraybackslash}p{0.158\linewidth} >{\raggedright\arraybackslash}p{0.158\linewidth} >{\raggedright\arraybackslash}p{0.158\linewidth} >{\raggedright\arraybackslash}p{0.158\linewidth} >{\raggedright\arraybackslash}p{0.158\linewidth} >{\raggedright\arraybackslash}p{0.158\linewidth}@{}}
\toprule
\textbf{Individual} & \textbf{age (kya)} & \textbf{TP53} & \textbf{CDKN2A} & \textbf{FOXO3} & \textbf{TERT} \\
\midrule
Present-day human (GRCh37) & 0 & 1.4333 & 1.4452 & 1.5927 & 1.5537 \\
Altai Neanderthal & \textasciitilde{}122 & 1.4325 & 1.4446 & 1.5913 & 1.5540 \\
Vindija 33.19 Neanderthal & \textasciitilde{}52 & 1.4325 & 1.4446 & 1.5928 & 1.5540 \\
Chagyrskaya 8 Neanderthal & \textasciitilde{}80 & – & – & – & 1.5540 \\
Denisova 3 (Denisovan) & \textasciitilde{}76 & 1.4327 & 1.4446 & 1.5921 & 1.5534 \\
Ust'-Ishim (early modern human) & \textasciitilde{}45 & 1.4333 & 1.4452 & 1.5921 & 1.5541 \\
Loschbour (W. hunter-gatherer) & \textasciitilde{}8 & 1.4333 & 1.4452 & 1.5927 & 1.5537 \\
\bottomrule
\end{longtable}
\end{center}

\subsection*{The senescence gate reads the same}

On TP53 (apoptosis) and CDKN2A (the p16 arrest switch), the two dated modern-human genomes and the present-day reference carry an identical promoter γ — the archaic individuals differ only by two or three distal substitutions that move γ by less than 0.0009. The core senescence/apoptosis gate is, as a measured state, the same sequence in every modern-human genome in the set and nearly the same in the archaic ones.

\subsection*{TERT carries the most promoter substitutions}

Of the four masters, the telomere-maintenance gene TERT carries the most homozygous-derived promoter substitutions across the archaic individuals (13 in total), and the same three TERT positions co-occur across all four archaic genomes. The widest γ range in the whole set, however, belongs to FOXO3 (0.001536), and even that is under 0.0016 — so “most substituted” (TERT) and “widest γ spread” (FOXO3) are distinct, both small, measured facts.

\begin{center}\footnotesize
\begin{longtable}{@{}>{\raggedright\arraybackslash}p{0.317\linewidth} >{\raggedright\arraybackslash}p{0.317\linewidth} >{\raggedright\arraybackslash}p{0.317\linewidth}@{}}
\toprule
\textbf{Gene} & \textbf{archaic hom-sub count} & \textbf{γ range (set)} \\
\midrule
TP53 & 8 & 0.000860 \\
CDKN2A & 6 & 0.000624 \\
FOXO3 & 6 & 0.001536 \\
TERT & 13 & 0.000712 \\
\bottomrule
\end{longtable}
\end{center}

\subsection*{The cancer-promoter positions are invariant}

The two recurrent TERT cancer-promoter regulatory positions (the −124 and −146 sites relative to the start codon) carry the same base in every individual in the set; every difference TERT does carry is distal to them. This invariance is a measured observation across the snapshot.

\begin{vpbox}\textbf{TERT is the most distinctive aging master} — on the cross-species axis (§9) it is the single residual lifespan lead; on this dated-individual axis it carries the most promoter substitutions — two independent measured observations, taken up in §11. [V] verified telomere keystone, §11\end{vpbox}

\subsection*{What is and is not claimed}

The claims are exactly three: the measured γ in each individual, the width of the γ spread per gene, and which promoter positions differ. Why the states co-occur is out of scope and graded [O]: a cross-sectional snapshot is silent on mechanism. Per-individual γ, the γ spread, the substitution counts, and the cancer-hotspot invariance are measured [V].

\section{The telomere keystone: dynamics, not gamma}

\begin{answerbox}The telomere is the keystone of aging on this substrate — but the keystone is its DYNAMICS, not its γ. The canonical telomere repeat (TTAGGG)n has γ 1.3298, the lowest of any aging-related sequence here and a near-exact universal constant (identical on both strands, length-independent to <0.5\%). TERT, which maintains this reservoir, is independently the most distinctive aging master on two measured axes (§9, §10), yet its promoter γ is a weak modulator because the telomere lever acts OFF the γ axis — through reservoir length and attrition rate.\end{answerbox}

\begin{abstractbox}γ fixes a node's barrier γ²/4 and reservoir capacity γ\textsuperscript{1.5} (DWELL). The telomere is a finite reservoir (§5); because its repeat γ is invariant, the reservoir RULER is the same everywhere and what differs between individuals and species is repeat COUNT (length) and attrition RATE, neither a γ quantity.\end{abstractbox}

This chapter is the package's synthesis of the question that motivated the telomere study: if the telomere is the key to aging, what exactly is the key — the γ of the machinery, or the dynamics of the reservoir it maintains?

\subsection*{The telomere repeat has the lowest, most invariant γ}

Run the promoter-γ rule on the canonical vertebrate telomere repeat (TTAGGG)n and the value is γ 1.3298. It is identical on the G-rich strand and its C-rich complement (the SantaLucia table is strand-symmetric), it is the lowest γ of any aging-related sequence in the package — below all four master promoters — and it is length-independent to better than half a percent (only end effects move it). The telomere γ is, in other words, a near-exact universal constant: the same ruler in every species and every dated individual.

\begin{center}\footnotesize
\begin{longtable}{@{}>{\raggedright\arraybackslash}p{0.475\linewidth} >{\raggedright\arraybackslash}p{0.475\linewidth}@{}}
\toprule
\textbf{Sequence} & \textbf{γ} \\
\midrule
\textbf{telomere repeat (TTAGGG)n} & 1.3298 \\
TP53 promoter & 1.4298 \\
CDKN2A promoter & 1.4424 \\
TERT promoter & 1.5539 \\
FOXO3 promoter & 1.5942 \\
\bottomrule
\end{longtable}
\end{center}

\subsection*{The reservoir reading: invariant ruler, variable length and rate}

On this substrate the telomere is a finite reservoir whose capacity scales as γ\textsuperscript{1.5} (DWELL) and which depletes monotonically to the replicative limit (§5). The run gives the telomere reservoir a capacity of 1.761 and 89 steps to depletion. But because the telomere-repeat γ is invariant, that γ\textsuperscript{1.5} ruler does not differ between individuals or species. What differs is the COUNT of repeats (telomere length) and the per-division attrition RATE — and neither of those is a γ or a sequence-identity quantity.

\begin{vpbox}\textbf{dwell(γ) ≈ γ\textsuperscript{1.5}} — how long a switch runs before its finite reservoir empties; it sets reservoir capacity (the replicative limit). [F] forced DNA Layer 1 (DOI)\end{vpbox}

\subsection*{TERT is the keystone gene — on two measured axes}

Two independent results in this package point at TERT. The cross-species discriminant (§9) finds it is the one master whose lifespan lean is not removed by body-mass correction — the single residual lead. The dated-individual observation (§10) finds it carries the most archaic promoter substitutions. On both axes TERT is the most distinctive aging master, and both are measured observations.

\subsection*{Yet the lever is off the γ axis}

Despite that, TERT's promoter γ sits inside the mammalian distribution (§9) and varies by less than 0.0008 across the dated-individual set (§10). There is no contradiction: the telomere lever on lifespan acts OFF the promoter-γ axis — through reservoir size (telomere length) and depletion rate (attrition) — the same off-axis pattern as the cross-species telomerase-dosage and TP53 copy-number levers (§9). The promoter-γ axis is therefore correct to report the telomere as only a weak γ modulator; the keystone role lives in the reservoir dynamics that γ does not index.

\begin{center}\ttfamily\small telomere γ = 1.3298 (lowest, universal)  |  keystone = reservoir dynamics, off the γ axis\end{center}

The telomere-repeat γ invariance, its strand symmetry, and its lowest-γ ranking are measured [V]; the γ\textsuperscript{1.5} reservoir law is the vendored substrate [F]; the “keystone is dynamics, not γ” synthesis organizes the package's own measured facts but is an interpretation [O]; telomere length and per-division attrition rates are cited, not reproduced in-package [L] (Harley 1990; Frenck 1998; Aubert \& Lansdorp 2008; somatic telomerase suppression in large mammals, Gomes 2011).

\section{Pathology: sarcopenia, frailty, cancer}

\begin{answerbox}The major non-rare aging diseases follow one derived law: g\_eff = γ(1−d) shrinks both the barrier (drift) and the spinodal (catastrophe). Sarcopenia is monotone setpoint drift, frailty is an accelerating co-failure of many setpoints, and aging-as-cancer is the convex crossing-hazard rise (135× from age 30 to 90). Rare/monogenic forms belong to disease\_wp.\end{answerbox}

\begin{abstractbox}A single derived law — an age-related loop-gain drop d lowering effective γ — generates three failure shapes: sarcopenia (defended force-state drifts down monotonically), frailty (critical drops spread across setpoints give an accelerating co-failure transition), and the convex 135× rise in crossing hazard that is aging-as-cancer-risk.\end{abstractbox}

\subsection*{One derived law}

Disease here is the failure of a defended setpoint on the R19 substrate. An age-related loop-gain drop d lowers the effective gain to γ(1−d); this shrinks the barrier (so the setpoint drifts) and the spinodal (so a fixed stressor can flip it). All three major failures are read off this one law.

\begin{vpbox}\textbf{barrier(γ) = γ\textsuperscript{2}/4} — the energy barrier between the two basins; it sets how stable a defended setpoint is. [F] forced substrate, §1\end{vpbox}

\begin{vpbox}\textbf{spinodal(γ) = 2(γ/3)\textsuperscript{1.5}} — the drive |h| past which one R19 basin disappears, so the flip is discontinuous (a saddle-node). [F] forced substrate, §1\end{vpbox}

\subsection*{Sarcopenia — monotone drift}

As gain falls, the defended force-state creeps down monotonically and the barrier fraction shrinks toward catastrophe. The shape is [V]; the absolute calendar rate is [O].

\begin{center}\footnotesize
\begin{longtable}{@{}>{\raggedright\arraybackslash}p{0.237\linewidth} >{\raggedright\arraybackslash}p{0.237\linewidth} >{\raggedright\arraybackslash}p{0.237\linewidth} >{\raggedright\arraybackslash}p{0.237\linewidth}@{}}
\toprule
\textbf{gain drop d} & \textbf{defended state} & \textbf{drift} & \textbf{barrier frac} \\
\midrule
0.0 & +1.1681 & +0.0000 & 1.000 \\
0.1 & +1.1019 & +0.0661 & 0.810 \\
0.2 & +1.0310 & +0.1370 & 0.640 \\
0.3 & +0.9540 & +0.2140 & 0.490 \\
0.4 & +0.8689 & +0.2992 & 0.360 \\
0.5 & +0.7716 & +0.3964 & 0.250 \\
0.6 & +0.6534 & +0.5146 & 0.160 \\
\bottomrule
\end{longtable}
\end{center}

\subsection*{Frailty and multimorbidity — accelerating co-failure}

Many setpoints decline together, each with its own critical drop (here spanning 0.28 to 0.58 across six systems). As the shared drop rises, the fraction of failed setpoints climbs in an accelerating S-curve — the frailty transition — rather than linearly. The co-failure shape is [V]; absolute prevalence is [O].

\subsection*{Aging as the cancer risk multiplier}

The crossing hazard rises convexly with age, a 135× increase from 30 to 90 in relative units, reproducing the steep oncology age-incidence slope (the seam of §7). The convex shape is [V]; the absolute incidence is [O].

Rare and monogenic accelerated-aging syndromes (progeroid disorders) are not modeled here; they are owned by the disease whitepaper and enter this layer only as a cited rate parameter.

\section{Reproducibility, grading, and the ledger}

\begin{answerbox}Every quantity is reproduced deterministically (SEED 19; the result hash is byte-identical across two runs at 62d5e1eb93db…) and graded honestly: mechanisms and the cross-species null result are [V], the copy-number longevity switch is [L], and absolute rates, incidence and lifespans are [O] with stated obstacles.\end{answerbox}

\begin{abstractbox}The discipline is LOCK → Derive → Gate with no tuning: every constant is a measured input or a derived value, never chosen to hit a target. The deterministic engine emits a byte-identical result on repeated runs (2×sha256), and the irreproducibility ledger collects every [O] item with its specific obstacle.\end{abstractbox}

\subsection*{LOCK → Derive → Gate}

Inputs are locked, results are derived by fixed rules, and gates check the derivation by count. The vendored substrate math is never re-derived here — the R19 field and its derived stability and dwell laws are defined once in §1; the promoter γ values are measurements; the dynamics are deterministic simulations seeded at 19.

\begin{center}\ttfamily\small SEED = 19  |  result sha256 = 62d5e1eb93db89630c4cd288… (2× identical)\end{center}

\subsection*{The grade vocabulary}

Three grades carry the evidence load. \textbf{[V]} is verified — measured or reproduced in-package. \textbf{[L]} is anchored to a cited result. \textbf{[O]} is open — not reproducible in-package, with the obstacle stated. The grading is the framework's load-bearing claim, not decoration.

\begin{itemize}[leftmargin=1.4em,itemsep=2pt,topsep=3pt]
\item \textbf{[V]} — the drift+catastrophe law (RA1), senescence irreversibility and accumulation (RA2), finite reservoirs and γ-ordering (RA3), the hallmark map (RA4), the convex risk-multiplier shape (RA5), the shared-rate structure (RA6), human-not-special and TP53 flatness (RA7), the archaic cross-sectional observation — per-individual γ, substitution counts, and TERT cancer-hotspot invariance (RA8) — and the telomere-repeat γ invariance, strand symmetry, and lowest-γ ranking (RA9).
\item \textbf{[L]} — the copy-number longevity switch (elephant \textasciitilde{}20 TP53 copies; Abegglen 2015 JAMA, Sulak 2016 eLife), the cited per-system decline anchors, and telomere length and per-division attrition rates (Harley 1990; Frenck 1998; Aubert \& Lansdorp 2008; somatic telomerase suppression in large mammals, Gomes 2011).
\item \textbf{[O]} — the γ–lifespan trend (small panel, GC-content confound, phylogenetic non-independence), the “telomere keystone is dynamics, not γ” synthesis (RA9, an interpretation that organizes the package's measured facts), the meaning of the archaic genotype co-occurrences (RA8, a cross-sectional snapshot is silent on mechanism), and every absolute calendar rate, incidence magnitude and lifespan, which need an external clock or calibration the package does not contain.
\end{itemize}

\subsection*{No tuning}

No constant is chosen to hit a target. Noise scales and chronic-stressor magnitudes used in the pathology sweeps are stated, round, and set the relative shape only; the absolute magnitudes they would imply are graded [O]. The irreproducibility ledger lists every [O] item, its obstacle, and its location.

\clearpage

\section{Conclusion: the genome fixes the ruler, not the lifespan}

\begin{answerbox}The genome fixes the ruler of aging — the promoter γ of the aging masters, the node inventory, and the reservoir law dwell ≈ γ\textsuperscript{1.5} — and that layer is measured from sequence [V]. It does not fix the realized lifespan: the rate at which senescent cells accumulate, the telomere attrition rate, and the absolute calendar age are runtime. The package's own results — γ nearly flat across a 4–211 yr lifespan span (TP53 CV 3.5\%), the telomere lever off the γ axis, and archaic promoters reading the same as modern — show that realized lifespan is underdetermined by the γ axis, not merely uncomputed; the longevity lever lives in copy number and reservoir dynamics, off the γ axis.\end{answerbox}

\begin{abstractbox}This closing section draws the package to its epistemic end. It separates two kinds of open item — \textit{uncomputed} (closable by an external clock or calibration) from \textit{underdetermined} (the γ axis does not carry the quantity) — and shows from the package's own measured results that the realized-lifespan gap is the second kind: across a fifty-fold cross-species lifespan span the flattest aging-master promoter γ varies by only CV 3.5\%, the telomere-repeat γ is a near-exact universal constant (1.3298) whose lever acts off the γ axis, and the four masters read the same γ in archaic and modern individuals. The longevity switch is copy number and reservoir dynamics, not promoter γ. The honest posture is observation under the constitution; an [O] is a correctly labelled truth, not a failure, and an open conclusion is the strongest position when the object is genuinely underdetermined.\end{abstractbox}

\begin{vpbox}\textbf{The reading of aging, taken to its honest end} — The genome fixes the ruler — the aging-master promoter γ, the node inventory, and the reservoir law dwell ≈ γ\textsuperscript{1.5} — all read from sequence [V]. The realized lifespan is the ruler run forward under boundary conditions the sequence does not contain: reservoir size (telomere length), attrition rate, copy-number dosage, and runtime drive. Two kinds of open item must not be confused: the absolute calendar rates are \textit{uncomputed} (they need an external clock), while the off-γ-axis location of the lifespan lever is \textit{underdetermined} — measured, and not closable by more computation on the promoter. The correct act is to observe and to grade honestly. [O] open reproducibility \&amp; grading, §13\end{vpbox}

\subsection*{γ is the ruler of aging, not its realized clock}

The aging masters' promoter γ is measured from sequence (§2), and it fixes each node's barrier γ²/4 and reservoir dwell ≈ γ\textsuperscript{1.5}. But γ does not index lifespan: across ten mammals spanning a fifty-fold maximum-lifespan range — from a mouse near 4 yr to a bowhead near 211 yr — the flattest master, TP53, varies at only CV 3.5\% (§9). A ruler that barely moves across a 50× lifespan span cannot be the axis that sets the lifespan. γ fixes the threshold scale, not the realized clock.

\subsection*{The telomere is the keystone of dynamics, not γ}

The single most distinctive aging master is TERT, on two measured axes: it is the residual cross-species lifespan lead (§9) and it carries the most archaic promoter substitutions (13, §10). Yet the canonical telomere repeat γ is 1.3298 — the lowest of any aging sequence and a near-exact universal constant, the same ruler in every species and every dated individual (§11). The telomere is the keystone of aging, but the keystone is its \textit{dynamics} — reservoir length and attrition rate — not its γ. The lever acts off the promoter-γ axis, exactly as the cross-species telomerase-dosage and TP53 copy-number levers do.

\subsection*{Archaic reads the same as modern — and that is silent on lifespan}

Across a cross-sectional set of seven dated genomes the four masters' promoter γ sits in a band narrower than 0.0016, and the senescence/apoptosis gate (TP53, CDKN2A) reads identically in the dated modern-human genomes (§10). This is a measured present-state, and it is silent on lifespan: realized lifespan was never on the γ axis, and the lifespans of archaic individuals cannot be observed at all. The honest reading is “the γ identity is the same,” not “the lifespan is the same.”

\subsection*{Two kinds of open, and the difference is the whole point}

An open item is not one thing, and conflating the two kinds makes the genome look either omniscient or irrelevant. An \textit{uncomputed} open item — the absolute calendar rates, incidences, and lifespans (§7, §8) — would close with an external clock or calibration the package deliberately does not contain; the shapes reproduce [V], but the absolute years do not. An \textit{underdetermined} open item is one the γ axis does not carry at all: the location of the lifespan lever off the γ axis is itself measured, and no more computation on the promoter sequence recovers it. And the cross-species null is not a broken test — it is audited robust to GC, body mass, phylogenetic non-independence, and multivariate confounds (§9), so the flatness is a result, not a failure to detect.

\subsection*{Neither all-genes nor environment}

The realized lifespan is therefore not “all in the genes” and not “the environment” either. The genome fixes the schedule, the node inventory, the thresholds, and even how open each node is to drive; the realized life is that ruler run forward under boundary conditions — reservoir size, attrition rate, copy-number dosage, and runtime drive — that the promoter sequence does not contain. The companion DNA reading reaches the same conclusion in the developmental-time axis: γ fixes what and in what order, not when.

\subsection*{Observation is the correct act, and an [O] is not a failure}

Where realization is underdetermined, the right act is to observe the measured states and grade honestly, not to manufacture an absolute lifespan the data withholds — which is exactly what the observation-only archaic comparison did. An [O] is not a failed [V]; it is a truth, correctly labelled, and it is what lets the framework learn from the next measurement rather than overclaim. An open conclusion is not a weak one: when the realized lifespan is genuinely underdetermined by the γ axis, it is the strongest position available, because it is the one the next measurement cannot overturn. The genome is the ruler and the first cause; the realized life is read off the ruler run forward under conditions the sequence does not hold. Not failure, not success: measurement.

\begin{center}\ttfamily\small genome fixes the ruler [V]  |  realized lifespan is off the γ axis [O] — dynamics, not γ\end{center}

\end{document}